\documentclass[11pt]{amsart}
\usepackage{amsmath,amssymb,amsthm}
\usepackage[margin=1.1in]{geometry}
\usepackage{booktabs}
\usepackage[colorlinks=true,linkcolor=blue,citecolor=blue,urlcolor=blue]{hyperref}
\hypersetup{pdftitle={Residual gap classes in the even-degree Fermat-fourfold census through degree 250}, pdfauthor={Rifat Jumagulov}}
\newtheorem{theorem}{Theorem}[section]
\newtheorem{lemma}[theorem]{Lemma}
\newtheorem{proposition}[theorem]{Proposition}
\newtheorem{corollary}[theorem]{Corollary}
\theoremstyle{remark}
\newtheorem{remark}[theorem]{Remark}
\newcommand{\Q}{\mathbb{Q}}
\newcommand{\Z}{\mathbb{Z}}
\newcommand{\F}{\mathbb{F}}
\newcommand{\PP}{\mathbb{P}}
\newcommand{\fr}[1]{\langle #1\rangle}
\newcommand{\V}{V}
\DeclareMathOperator{\Gal}{Gal}
\newcommand{\zt}{\zeta}
\newcommand{\C}{\mathbb{C}}
\newcommand{\Qbar}{\overline{\Q}}
\vfuzz=2pt

\title[Even-degree Fermat gap classes]{Residual gap classes in the even-degree Fermat-fourfold
census through degree 250: exchange walls and two closures}
\author{Rifat Jumagulov}
\email{jum.rifm@gmail.com}
\date{July 2026}

\begin{document}
\raggedbottom
\begin{abstract}
For the Fermat fourfold $X^4_m\colon x_0^m+\dots+x_5^m=0$, the companion note \cite{BN} proves
that the sector beyond the classical algebraicity machinery is empty in odd degree $\le199$.
This paper is the systematic treatment of the \emph{even}-degree sector through degree $250$,
which behaves differently: gap classes are frequent, and the odd law ``survivor
$\Rightarrow3\mid m$'' fails. We first identify the working lattice with Aoki's group $S_m$ generator-by-generator ---
including his $p=2$ standard elements, which we show are identically grade-$2$ Hodge quadruples,
hence Lefschetz-algebraic --- and prove a negation-closure lemma converting any integer
membership certificate into a literal Aoki claim chain. The even-parity census --- run under the
note's completeness lemma, which applies to both parities; receipts and independent verifiers in
the ancillary --- finds
thirty-eight primitive beyond-machinery classes through degree $250$; iterated partition
certificates close all but ten. An exchange theorem generates claim-equivalences, and an
exhaustive scan organizes the ten into five walls. Two of the five then close: the $m=168$ class
by Aoki's published degree-form theorem (citing which we record two errata of the printed text,
\S\ref{sec:closures}), and the isolated
degree-$210$ class by an explicit $40$-generator $\pm1$ certificate in $S_{210}$, displayed in
full. The residual even-degree boundary at $m\le250$ --- the classes unresolved by the
implemented closure family and the cited published theorems --- is eight certified gap classes
($\nu\notin S_m$,
$2\nu\in S_m$) in three exchange walls $W_{70}$, $W_{110}$, $W_{114}$; every lattice verdict
carries an independent certificate (modular kernel witnesses for non-membership, exact
re-summation for membership).
\end{abstract}
\maketitle

\section{Introduction}\label{sec:intro}

Let $X^4_m\subset\PP^5$ be the Fermat fourfold, $G=\mu_m^6/\mu_m$, and for a character
$a=(a_0,\dots,a_5)\in(\Z/m)^6$ with $\sum a_i\equiv0$, $a_i\neq0$, let
$\V(a)\subset H^4_{\mathrm{prim}}$ be the one-dimensional eigenspace; $a$ is a Hodge $(2,2)$
character iff $\sum_i\fr{ta_i}=3m$ for all $t\in(\Z/m)^\times$ (Shioda \cite{Shioda79}). The
classical algebraicity machinery --- Shioda's inductive structure with da Silva's
quasi-decomposability formalism \cite{Shioda79,SK,daSilva} and Aoki's standard cycles
\cite{Aoki87} --- and the notation of this paper are those of the companion note \cite{BN}, whose
conventions we use throughout: characters are \emph{ordered} tuples; $\mathrm{claim}(a)$
($\V(a)$ spanned by algebraic cycle classes) is invariant under coordinate permutations
(automorphisms of $X^4_m$) and is a Galois-orbit-level property; the census classifies sorted
representatives modulo both actions; \emph{grade} $g$ means $\sum_i\fr{ta_i}=gm$ for every unit
$t$; $*$ is juxtaposition of entry tuples, $\uplus$ multiset union; $\chi_a$ is the Jacobi-sum
Hecke character of $a$ and $u(a,p)$ its Tate-normalized value at a split prime $p$
\cite{Weil49,Weil52}, while $\nu(a)\in\Z^{m-1}$ is Aoki's multiplicity vector, additive under
juxtaposition (he writes $u$ for $\nu$); and the inflation/descent lemmas of
\cite{BN} make claim level-transparent along $x\mapsto x^e$. The note proves (Theorems A, A$'$,
A$''$, B there) that every Hodge $(2,2)$ class of every $X^4_m$ of \emph{odd} degree $m\le199$
is algebraic.

This paper maps the even-degree sector through degree $250$ and proves its structural results.
The even sector is a genuinely different landscape: the first beyond-machinery classes appear at
$m=50$ (so ``survivor $\Rightarrow3\mid m$'', true throughout the odd census, fails), most
beyond-machinery classes close by \emph{deeper} partition certificates rather than new
mechanisms, and what finally remains is organized not by degree but by \emph{claim-equivalence}.
Our results:

\begin{itemize}
\item[--] \emph{The lattice is Aoki's} (Proposition~\ref{prop:ident}): the working lattice ---
pair vectors and standard vectors, the $p=2$ family included --- coincides
generator-by-generator with the group $S_m$ of \cite[\S2]{Aoki00}, so all membership and gap
verdicts are statements about Aoki's own object. The $p=2$ standard elements are identically
grade-$2$ Hodge quadruples, hence Lefschetz-algebraic (Lemma~\ref{lem:p2neg}(i)), and the
standard family is negation-closed modulo pairs (Lemma~\ref{lem:p2neg}(ii)).
\item[--] \emph{The census} (Proposition~\ref{prop:evencensus}): thirty-eight primitive
beyond-machinery classes through degree $250$; iterated depth-$3/4$ partition certificates close
all but ten.
\item[--] \emph{Exchange walls} (Theorem~\ref{thm:exchange}, Proposition~\ref{prop:walls}): a
two-entry (or three-entry) exchange with a Lefschetz-algebraic exchanger is a claim-equivalence;
the exhaustive scan organizes the ten residual classes into five walls.
\item[--] \emph{Two closures} (Propositions~\ref{prop:w168} and~\ref{prop:w210}): the $m=168$
class falls under Aoki's published Theorem~0.1(i) \cite{Aoki00} --- we record two errata of
the printed text, with a machine-verified possible replacement for the $m=28$
gap-generator misprint --- and the isolated degree-$210$ class closes by an explicit
$40$-generator $\pm1$ lattice certificate, displayed in full and converted to a literal Aoki
chain.
\item[--] \emph{The boundary} (Corollary~\ref{cor:evenbdry}): eight certified gap classes in
three walls; closing any member of a wall closes the wall.
\end{itemize}

Status ledger: the Propositions labeled \emph{computer-assisted} are exhaustive machine
computations with pinned receipts, engines, and independent verifiers in the ancillary files;
their scope (which levels, which move family) is stated in each. Everything else is proved.
The obstruction theory of the closed odd classes (exact Frobenius certificates, the negative
answer to da Silva's Question~1 at $m=33$) is in \cite{BN} and not repeated here.

\section{The lattice is Aoki's $S_m$}\label{sec:lattice}

$S_m$ denotes Aoki's lattice: the subgroup of $\Z^{m-1}$ generated by the $\nu(\delta)$ over the
vanishing pairs $\delta=\{k,m-k\}$ (the self-pair included at even $m$) and the
$\nu(\sigma_{p,x})$ over the standard characters, for all primes $p\mid m$; membership is decided
exactly by integer linear algebra (Hermite normal form), and in lattice
statements we freely identify $a$ with $\nu(a)$ --- statements about the formal calculus, not
about the character group. The first result pins the object.

\begin{proposition}[the implemented lattice is Aoki's $S_m$]\label{prop:ident}
Aoki \cite[\S2]{Aoki00} defines $A_m$ as the zero-sum sublattice of the free abelian group on
$\Z/m\smallsetminus\{0\}$; $u$ --- our $\nu$ --- as the multiplicity-vector map, additive
under juxtaposition, $\nu(\alpha*\beta)=\nu(\alpha)+\nu(\beta)$; the standard elements,
verbatim, as
$\sigma_{p,a}=(a,\,a+\frac mp,\dots,a+\frac{(p-1)m}p,\,m-pa)$ for $p\ge3$ and
$\sigma_{2,a}=(a,\,a+\frac m2,\,m-2a,\,\frac m2)$ for $p=2$, over primes $p\mid m$ ($m>p$) and
$0<a<\frac mp$ with nonzero entries; $\mathfrak S_m$ as the set of $\alpha$ admitting
$\alpha*\delta\sim\sigma_1*\cdots*\sigma_k*\delta'$ with $\delta,\delta'$ unions of vanishing
pairs, with $\mathfrak D_m\subset\mathfrak S_m$; and $S_m\subset A_m$ as the subgroup generated
by $\nu(\mathfrak S_m)$. Consequently $S_m$ equals the $\Z$-span of the pair vectors
$\nu(\{k,m-k\})$ and the standard vectors $\nu(\sigma_{p,a})$, $p=2$ included --- exactly the
generator list of the implemented lattice --- so the machine verdicts $\nu\in S_m$ /
$\nu\notin S_m$ are statements about Aoki's group itself.
\end{proposition}
\begin{proof}
$\supseteq$: every pair class lies in $\mathfrak D_m\subset\mathfrak S_m$, and every standard
element $\sigma$ satisfies $\sigma*\delta\sim\sigma*\delta$, so $\sigma\in\mathfrak S_m$: each
listed generator is in $\nu(\mathfrak S_m)$. $\subseteq$: if
$\alpha*\delta\sim\sigma_1*\cdots*\sigma_k*\delta'$, then
$\nu(\alpha)=\sum_i\nu(\sigma_i)+\nu(\delta')-\nu(\delta)$, an integer combination of the listed
generators. The two spans coincide.
\end{proof}

\begin{lemma}[$p=2$ standards; negation closure]\label{lem:p2neg}
\begin{enumerate}
\item[(i)] Every $\sigma_{2,x}$ is a grade-$2$ Hodge quadruple identically in $x$, so
$\V(\sigma_{2,x})$ is algebraic by Lefschetz $(1,1)$ on the Fermat surface $X^2_m$: the $p=2$
generators carry the same unconditional algebraicity as the odd-$p$ standards (Theorem~2-1 of
\cite{Aoki87}, with the inflation lemma of \cite{BN} when $\gcd(x,m/p)>1$).
\item[(ii)] For every prime $p\mid m$ and admissible $x$, $\sigma_{p,x}\uplus\sigma_{p,\,m/p-x}$
is a union of vanishing pairs --- $p+1$ of them for odd $p$, four for $p=2$; in the lattice,
$-\nu(\sigma_{p,x})\equiv \nu(\sigma_{p,\,m/p-x})\bmod D_m$, so an integer membership certificate
converts to one with nonnegative standard coefficients at the cost of pair classes.
\end{enumerate}
\end{lemma}
\begin{proof}
(i) For odd $t$ write $r=\fr{tx}$; $t\cdot\frac m2\equiv\frac m2$, and the residues of
$t\cdot\sigma_{2,x}$ are $r,\,r+\frac m2,\,m-2r,\,\frac m2$ if $r<\frac m2$ and
$r,\,r-\frac m2,\,2m-2r,\,\frac m2$ if $r>\frac m2$ ($r=\frac m2$ is excluded: $x+\frac m2
\not\equiv0$) --- both sum to $2m$. Grade-$2$ Hodge quadruples are the $(1,1)$ characters of
$X^2_m$; the rational part of their Galois block is divisorial by Lefschetz $(1,1)$, whence the
claim. (ii) For odd $p$, $d=m/p$: the translates pair off as
$(x+jd)+\bigl((d-x)+(p-1-j)d\bigr)=pd=m$ for $j=0,\dots,p-1$, and $(m-px)+(px)=m$. For $p=2$:
$\{x,m-x\}$, $\{x+\frac m2,\frac m2-x\}$, $\{m-2x,2x\}$, and the self-pair
$\{\frac m2,\frac m2\}$. Each pair lies in the span of the pair vectors.
\end{proof}

\begin{remark}[even-$m$ standard admissibility]\label{rem:evenstd}
For even $m$ and odd $p\mid m$ Aoki's 1987 admissibility hypothesis $d/\gcd(x,d)>2$ ($d=m/p$) is
strictly stronger than ``entries nonzero'': it excludes $x\equiv d/2\pmod d$. Every excluded
multiset contains a vanishing pair --- $(x+k_1d)+(x+k_2d)\equiv0\pmod m$ reduces to
$1+2j+k_1+k_2\equiv0\pmod p$ for $x=d/2+jd$, solvable with $0\le k_1<k_2\le p-1$ --- so it is
decomposable, hence algebraic for the trivial reason, and the census classifier (which tests
decomposability first) is unaffected. The lattice of Proposition~\ref{prop:ident} follows the
range of \cite[\S2]{Aoki00} (nonzero entries only), matching his definition of $S_m$.
\end{remark}

\begin{remark}[the gap group]\label{rem:gap}
Aoki's Theorem~2.1 \cite{Aoki00} (after \cite[Thm.~D]{AokiD}; the gap-group structure goes back
to Yamamoto \cite{Yamamoto}) proves $2\alpha\in S_m$ for \emph{every} Hodge class $\alpha\in
B_m$: the gap group $B_m/S_m$ is an elementary abelian $2$-group, generated by level-raised gap
generators $\xi_d$. Every ``$2u\in S_m$'' certificate below therefore instantiates a general theorem; the
content of a \emph{gap} verdict is the non-membership half $u\notin S_m$.
\end{remark}

\section{The even census}\label{sec:census}

For each even $m$: enumerate grade-$3$ Hodge multisets (meet-in-the-middle on half-unit sum
profiles); reduce to Galois orbits of sorted representatives; classify
decomposable $\to$ quasi-decomposable $\to$ standard $\to$ $*$-split (the transport theorem of
\cite{BN}, stated there for all $m$) $\to$ survivors. The
census-completeness lemma of \cite{BN} is stated and proved there for both parities (the even
self-pair $m/2$ included in $M_m(1)$); the even-parity engine is calibrated against direct brute
force --- exact match of representative counts for $m\le14$ --- while the independent blind
reimplementations of \cite{BN} cover the odd census.
Throughout, \emph{primitive} means content $\gcd(a_0,\dots,a_5,m)=1$, i.e.\ not an inflation
from a lower level; eigenvalue-field labels are \emph{numerical observations} (an mpmath
heuristic over finitely many split primes), used as descriptors only --- no formal statement
rests on them. Two engine-tier facts are recorded for transparency. The receipt engine (the v3
pipeline over the even-safe base library) is multiplicity-aware throughout: its decomposability
requires both members of a vanishing pair, the self-paired $m/2$ requiring multiplicity two, and
its quasi loop includes the self-pair $\{m/2,m/2\}$ (e.g.\ $(1,3,36,38,64,68)$ at $m=70$ is
genuinely quasi-decomposable through the self-pair, and the engine closes it). The odd-heritage
display engine and the first version of the independent reimplementation carried set-based tests
unsound at even $m$ --- a single $m/2$ reads as a vanishing pair; regression class
$(1,29,35,43,45,57)$ at $m=70$, neither decomposable nor quasi-decomposable. The independent
reimplementation is fixed (multiplicity-aware), both regression classes are asserted afresh by
the runner, and its fresh key-level runs reproduce the engine receipt's representative and
base-closure counts exactly at $50,70,110,114,168$ (pinned receipt). Independently of any
engine, an exhaustive self-pair-inclusive quasi test --- every pair, every split, all units ---
certifies all ten residual-table classes non-quasi (pinned receipt): no wall is affected. For
self-containedness we recall the imported transport statement: \emph{a Hodge $(2,2)$ character
of $X^4_m$, $m$ arbitrary, whose $6$-multiset splits into two zero-sum triples has algebraic
rational Hodge classes in its Galois block} (\cite[Thm.~A]{BN}; its transport uses only the
Galois-stable block, Lefschetz $(1,1)$, and the $\Q$-linear joining-line correspondence, so no
rational point on the curve and no parity hypothesis enters). Counting convention: \emph{beyond-machinery} counts the primitive
orbits surviving all four base predicates (decomposable / quasi-decomposable / standard /
$*$-split), before the iterated and deep tiers; the paper's $38=21+14+3$ is this base-tier
primitive count over $m\le108$, $108<m\le200$, $200<m\le250$. Surviving orbits are then attacked by the \emph{closure
operator}: $T$ closes at depth $k$ if
$T\uplus(k\text{ vanishing pairs})$, $k\le4$, partitions into blocks each of which is a
vanishing pair, a grade-$2$ Hodge quadruple (Lefschetz $(1,1)$), an admissible standard block
(Theorem~2-1 of \cite{Aoki87}; Lemma~\ref{lem:p2neg}(i) at $p=2$), or a previously certified
sextuple (iteration: two-generation certificate trees); $\mathrm{claim}(T)$ then follows by
Theorems~1-4(i),(ii) of \cite{Aoki87}. Every closure is stored as an explicit route in the
deep-closure ledger and re-verified by the standalone checkers.

\begin{proposition}[even census]\label{prop:evencensus}
{\rm(Computer-assisted; engines, closure ledger, and receipts in the ancillary.)}
(i) For even $m\le108$ there are twenty-one primitive beyond-machinery orbits, beginning at
$m=50$, with indicated eigenvalue fields up to degree $32$; the odd law ``survivor
$\Rightarrow3\mid m$'' fails in the even sector. (ii) Eighteen of the twenty-one close by
depth-$3/4$ partition certificates (two closed sextuples, or quadruple--standard--quadruple);
iterating the closure operator (blocks certified by their own certificates --- two-generation
trees) closes $m=54$ and $m=90$ as well, so exactly one even class of degree $\le108$ resists:
$m=70$, $(1,20,24,42,61,62)$, indicated eigenvalue field $\Q(\zt_{35})$. (iii) The census through degree
$200$ yields fourteen further primitive classes, of which the same machinery closes eight (every
open row in $178$--$200$ is an induced copy of an already-closed class); the residual at
$m\le200$ is seven classes ($70$; two at $110$; three at $114$; one at $168$) and their induced
copies. (iv) Through degree $250$ exactly three more primitive classes appear:
$(2,9,129,142,168,180)$ and $(1,79,109,121,151,169)$ at $210$, and $(1,62,111,142,162,182)$ at
$220$.
\end{proposition}

Among the eighteen closures of (ii) is the $m=66$, $\Q(\sqrt{-3})$ class --- closed through the
self-paired element $(33,33)$, the same resource that closes da Silva's $m=33$ witness at level
$66$ in \cite[Thm.~A$''$]{BN} --- so at that census depth every known Fermat-fourfold Hodge
class with imaginary-quadratic \emph{indicated} eigenvalue field was algebraic; the deeper
census produces a third such class at $m=168$, closed in \S\ref{sec:closures}. That survivor,
$(1,25,79,121,127,151)$, has indicated eigenvalue field $\Q(\sqrt{-3})$, and its Frobenius
character has a kernel prime ($p=673$), so the field-of-definition obstruction of
\cite[Thm.~C]{BN} vanishes there.

\section{Exchange and the wall structure}\label{sec:exchange}

\begin{theorem}[Exchange]\label{thm:exchange}
Let $T=S\uplus A$ and $T'=S\uplus B$ be Hodge sextuples $\bmod\ m$ (entries nonzero) with
$|A|=|B|=2$, and suppose $q:=A\uplus(-B)$ has nonzero entries and is a Hodge $(1,1)$ quadruple.
Then $\mathrm{claim}(T')\Rightarrow\mathrm{claim}(T)$; since $-q$ serves the reverse direction,
$T$ and $T'$ are \emph{claim-equivalent}. The same holds for $|A|=3$ with $q$ a sextuple whose
claim is known.
\end{theorem}
\begin{proof}
$V(q)$ is algebraic by Lefschetz $(1,1)$ on the Fermat surface $X^2_m$ (no membership in
$\mathfrak D_m$ is required). By Theorem~1-4(i) of \cite{Aoki87} (juxtaposition join, $r,s$ even),
$\mathrm{claim}(T')$ and $\mathrm{claim}(q)$ give $\mathrm{claim}(T'\ast q)$. As multisets,
$T'\uplus q = T\uplus(y_1,-y_1)\uplus(y_2,-y_2) = T\ast\delta$ with
$\delta=(y_1,-y_1,y_2,-y_2)$ decomposable, so $\delta\in\mathfrak D^2_m$ and Theorem~1-4(ii) of
\cite{Aoki87} cancels it. For $|A|=3$ the same two steps apply with $q$ a sextuple of known claim
and $\delta=B\uplus(-B)$ three vanishing pairs, $\delta\in\mathfrak D^4_m$ ($s=4$ even); and
$\mathrm{claim}(-q)=\mathrm{claim}(q)$ since $-q=(-1)\cdot q$ is a Galois conjugate. (Landing on a
Galois conjugate $t\cdot T'$ suffices, claim being orbit-level; the
$(m/2,m/2)$ self-pair is a legitimate multiset entry.)
\end{proof}

\begin{proposition}[wall structure through degree 250]\label{prop:walls}
{\rm(Computer-assisted: exhaustive exchange scan, both $|A|$ shapes, own levels and level-lifts.)}
The two $m=110$ classes are claim-equivalent --- the $\Q(\zt_5)$ class stands or falls with
the $\Q(\zt_{110})$ one --- and the three $m=114$ classes form a complete triangle at
Galois-orbit level, while $m=70$ and $m=168$ are exchange-isolated. The three primitive classes
of Proposition~\ref{prop:evencensus}(iv) are already placed by the exchange graph:
$(2,9,129,142,168,180)$ at $210$ is equivalent to the $3$-lift of the $m=70$ class; the $220$
class is equivalent to the $2$-lifts of the $m=110$ pair; $(1,79,109,121,151,169)$ at $210$ is
new and isolated. The residual boundary after the census machinery --- before the closures of
Propositions~\ref{prop:w168} and~\ref{prop:w210} --- is thus ten primitive classes organized in
five claim-equivalence walls
\begin{gather*}
W_{70}=\{70\simeq210\#2\},\qquad W_{110}=\{110\#1\simeq110\#2\simeq220\},\\
W_{114}=\{\#1\simeq\#2\simeq\#3\},\qquad W_{168},\qquad W_{210}=\{210\#1\},
\end{gather*}
and closing any member of a wall closes the wall (Theorem~\ref{thm:exchange}). Precisely: call
$T\leftrightarrow T'$ an \emph{exchange move} when $T=S\uplus A$, $T'=S\uplus B$ with
$|A|=|B|\in\{2,3\}$ and the exchanger $q=A\uplus(-B)$ has nonzero entries and is a grade-$2$
Hodge quadruple ($|A|=2$; unconditional by Lefschetz) or a grade-$3$ sextuple of known claim
($|A|=3$). The scan enumerated \emph{every} move out of \emph{every} open class --- all entry
positions, all $B$ --- at its own level and its lifts within the census range (pinned receipts:
the own levels $70,110,114,168,210$ and the lifted levels $140,220,228$); the open set supplied
to the scan at each level is the complete final list of unresolved orbits there, so an $|A|=3$
exchanger outside it carries a known algebraicity claim; each lands on the
class's own wall, never on a closed class. Within this move family and level range, the
transitive exchange closure of each open class is therefore exactly its wall.
\end{proposition}

\section{Two closures}\label{sec:closures}

\begin{proposition}[the $m=168$ class closes]\label{prop:w168}
The class $(1,25,79,121,127,151)$ on $X^4_{168}$ is algebraic; with it every class of the two
census ranges (odd $m\le199$, even $m\le250$) whose indicated eigenvalue field is imaginary
quadratic ($m=33,66,168$) is closed.
\end{proposition}

\begin{proof}
The closure rests on the published theorem: Aoki \cite{Aoki00} proves (Theorem~0.1(i)) the Hodge
conjecture for all abelian varieties of Fermat type of degree $m=2^a3^b5^c7^d$ ($c=0$ or $d=0$),
and the proof's \S4 route opens with the bridge through Corollary~3.2 there, establishing the
statement for the Fermat varieties $X^n_m$ themselves, for all $n$, at these degrees; $168 =
2^3\cdot3\cdot7$ qualifies. We do not re-derive that proof; we record two errata of the printed
text, both verified against the primary and neither affecting the theorem as stated. First, the
paper's own later citation of this statement
(\cite[Thm.~1.2(ii)]{Aoki02}, attributed to a then-preprint on Catalan curves) is a
citation slip --- the published Catalan paper (Tokyo J.\ Math.\ \textbf{27} (2004)) concerns the
Jacobians of generalized Catalan curves, not the degree form; the theorem is \cite{Aoki00}'s own
Thm.~0.1(i). Second, in the \S4 list of gap generators the $m=28$ entry
$\xi_{28}=(1,9,25,12,20,24)$ is a misprint (it is not a zero-sum character), while Aoki's own
text requires $\xi_m\in B^2\cup\bigl(B^4\cap(\mathfrak A^1*\mathfrak A^1)\bigr)$ and algebraizes
these $\V(\xi_m)$ by the proof of Theorem~4.3 of Shioda \cite{Shioda81} (cited there); a
machine-verified \emph{possible replacement} of exactly the required join form is
$\xi_{28}=(1,9,18)\ast(10,21,25)\in B^4_{28}\cap(\mathfrak A^1_{28}*\mathfrak A^1_{28})$,
$\nu(\xi_{28})\notin S_{28}$, $2\nu(\xi_{28})\in S_{28}$ (whether it is the intended
generator modulo $S_{28}$ is not established here; the closure of $168$ does not depend on
it). At the class level our lattice computation
gives $\nu(a)\notin S_{168}$ with $\nu(a)\equiv \nu(6\cdot\xi_{28})\bmod S_{168}$ (entrywise product
--- the level-raising map of \cite[\S2]{Aoki00}): the class is a genuine gap class, and no
Fermat-internal route of this paper or of \cite{BN} reaches it (its multiset admits no
zero-sum-triple split, so the $*$-split theorem --- stated in \cite{BN} for all $m$ --- does not
fire, machine-checked; $5\nmid168$ leaves no standard block for the two-pair genre; the
iterated and lift searches also failed) --- Theorem~0.1(i) is what closes it.
\end{proof}

\begin{proposition}[the isolated degree-210 class closes]\label{prop:w210}
$\V(1,79,109,121,151,169)$ on $X^4_{210}$ is algebraic.
\end{proposition}

\begin{proof}
The class closes by an explicit
certificate: $\nu(1,79,109,121,151,169) \in S_{210}$, realized by $40$ generators with coefficients
$\pm1$ (re-summed independently in the ancillary; stored as a data artifact). In full, with $(k)$
denoting the pair $\{k,210-k\}$:
\begin{gather*}
\scriptstyle \nu(a)\;=\;
-(7)-(8)+(12)-(21)-(27)-(28)-(29)+(36)+(41)+(42)-(61)\\
\scriptstyle
+(66)+(74)-(76)-(77)-(78)+(82)+(84)-(91)+(96)-(105)\\
\scriptstyle
+\,\sigma_{2,4}-\sigma_{2,9}+\sigma_{2,14}+\sigma_{2,21}+\sigma_{2,27}+\sigma_{2,28}+\sigma_{2,29}
-\sigma_{2,31}-\sigma_{2,41}-\sigma_{2,42}\\
\scriptstyle
-\sigma_{3,4}+\sigma_{3,6}+\sigma_{3,7}+\sigma_{3,8}+\sigma_{3,9}-\sigma_{3,12}-\sigma_{3,14}
+\sigma_{7,1}-\sigma_{7,6}\,.
\end{gather*}
Every generator's eigenspace is algebraic --- pairs: linear cycles, Theorem~1-1 of \cite{Aoki87};
$\sigma_{3,x},\sigma_{7,x}$: Theorem~2-1 of \cite{Aoki87}, with the inflation lemma of \cite{BN}
when $\gcd(x,m/p)>1$; $\sigma_{2,x}$: Lemma~\ref{lem:p2neg}(i) --- so converting the negative
terms by Lemma~\ref{lem:p2neg}(ii) produces the literal
membership $a\ast\delta \sim \sigma_1\ast\cdots\ast\sigma_{19}\ast\delta'$ ($\delta,\delta'$
unions of vanishing pairs), and $\V(a)$ is algebraic by Theorems~1-4(i),(ii) of \cite{Aoki87} ---
the calculus of Lemma~4.1 of \cite{Aoki00} (which carries no hypothesis on $m$), with every block
certified above.
\end{proof}

Negative-coefficient certificates of this kind are invisible to
monoid/partition search engines --- the lattice membership test ($\nu(a)\in S_m$, Hermite normal
form) decides standard-calculus reachability exactly and is the instrument we now run first. (The
two directions, with the scope stated precisely: a certificate all of whose blocks are vanishing
pairs, standard blocks, or classes already certified inside $S_m$ forces $\nu(a)\in S_m$, since $\nu$
is additive under $*$; a certificate routed through a Lefschetz quadruple whose own $\nu$ lies
outside $S_m$ escapes the lattice --- the $\Q(\zt_{21})$ class at $105$ closes exactly this way
in \cite[Thm.~A$'$]{BN} while remaining a gap class --- so gap-ness bounds the \emph{standard}
calculus, not partition search as such. Conversely an integer membership combination converts to
a literal Aoki chain by negation closure (Lemma~\ref{lem:p2neg}(ii)), as carried out in
Proposition~\ref{prop:w210}.)

\section{The boundary at $m\le250$}\label{sec:boundary}

\begin{corollary}[the even-degree boundary at $m\le250$]\label{cor:evenbdry}
The residual boundary for all \emph{even} degrees $m\le250$ --- unresolved by the implemented
closure family and the cited published theorems --- is eight primitive classes in three
walls ($W_{70}$: two members, $W_{110}$: three, $W_{114}$: three, including the displayed
lift-equivalent members at $210$ and $220$), and
all eight are certified gap classes: $\nu\notin S_m$, $2\nu\in S_m$ (the doubling holds for
\emph{every} Hodge class by \cite[Thm.~2.1]{Aoki00}; the certificates make it explicit). The
residual question is uniform: algebraicity of gap-group
classes at degrees $70, 110, 114$ --- outside the degree forms of Aoki's theorem ($70$ and $210$
contain both $5$ and $7$; $110$ contains $11$; $114$ contains $19$).
\end{corollary}

\medskip
\begin{center}\small
\begin{tabular}{llll}
\toprule
$m$ & class & $\nu\in S_m$? & status\\
\midrule
$70$  & $(1,20,24,42,61,62)$      & no  & unresolved ($W_{70}$)\\
$110$ & $(1,24,62,71,81,91)$      & no  & unresolved ($W_{110}$)\\
$110$ & $(1,31,55,71,81,91)$      & no  & unresolved ($W_{110}$)\\
$114$ & $(1,13,43,72,103,110)$    & no  & unresolved ($W_{114}$)\\
$114$ & $(1,13,43,80,102,103)$    & no  & unresolved ($W_{114}$)\\
$114$ & $(1,7,78,79,86,91)$       & no  & unresolved ($W_{114}$)\\
$168$ & $(1,25,79,121,127,151)$   & no  & closed (Prop.~\ref{prop:w168})\\
$210$ & $(1,79,109,121,151,169)$  & yes & closed (Prop.~\ref{prop:w210})\\
$210$ & $(2,9,129,142,168,180)$   & no  & unresolved ($W_{70}$)\\
$220$ & $(1,62,111,142,162,182)$  & no  & unresolved ($W_{110}$)\\
\bottomrule
\end{tabular}

\smallskip
\parbox{0.92\linewidth}{\footnotesize The ten primitive classes of degree $\le250$ beyond the
base machinery. Every lattice verdict carries an independent certificate in the ancillary: a
modular kernel witness ($q=2$) for each ``no'', the re-summed $\pm1$ combination for the ``yes''.}
\end{center}

A systematic treatment of the even sector beyond degree $250$ is deferred.

\section{Methods, code availability, and provenance}\label{sec:methods}
Exact objects throughout: multiset identities, Hodge grades, and lattice membership are decided
in exact integer arithmetic; no floating point enters any certificate. Every lattice verdict has an
independent exact certificate (modular kernel witnesses found by the standalone script's own
elimination for each non-membership; exact re-summation for each membership and for the doubled
vectors); the full census and exchange sweeps are reproducible from the supplied engines and are
accompanied by checksummed campaign receipts; independent brute-force calibration is performed
at small levels ($m\le14$), and the algorithmically independent reimplementation re-derives the
representative counts fresh at the key levels $50,70,110,114,168$ (pinned receipt). The author
takes full responsibility for all proofs, code, and bibliographic claims. The ancillary package is shared with the companion note: engines (census, closure
oracle, exchange scan, lift sweeps, lattice membership), the deep-closure ledger, the stored
degree-$210$ certificate, the complete even census receipt over all even $6\le m\le250$, and the
pinned exchange/negative-sweep logs, with a checksum manifest and a fresh-run runner
(\texttt{run\_all.sh}; ``PASSED'' always means executed now, never cached). This manuscript was
prepared with AI assistance; all computational claims are backed by the ancillary code and
receipts.

\appendix
\section{The cited Aoki statements}\label{app:aoki}

For self-containedness we reproduce the exact statements used, in the numbering of the primary
texts.

\emph{From \cite{Aoki87}} (notation as there: $\mathfrak B^n_m$ the Hodge characters of $X^n_m$,
$\mathfrak D^n_m$ the classes that are coordinate permutations of $(a_0,-a_0,\dots,a_r,-a_r)$,
$r=n/2$; $*$ juxtaposition; \textsc{claim}$(\alpha)$: $\V(\alpha)$ is generated by classes of
algebraic cycles on $X^n_m$):

\begin{itemize}
\item[--] \textbf{Theorem 1-1} (Shioda). \emph{The linear space $L$ represents
$\delta\in\mathfrak D^n_m$; more precisely
$\omega_\delta(L)\cdot\overline{\omega_\delta(L)}=(-1)^r m^{n+1}$.}
\item[--] \textbf{Theorem 1-4} (Shioda, Ran). \emph{Let $r$ and $s$ be non-negative even integers
such that $n=r+s+2$, and let $\alpha\in\mathfrak B^r_m$, $\beta\in\mathfrak B^s_m$. Then:
(i) if \textsc{claim}$(\alpha)$ and \textsc{claim}$(\beta)$ are true, then
\textsc{claim}$(\alpha*\beta)$ is also true; (ii) if there exists $\delta\in\mathfrak D^s_m$ such
that \textsc{claim}$(\alpha*\delta)$ is true, then \textsc{claim}$(\alpha)$ is also true.}
\item[--] \textbf{Theorem 2-1.} \emph{For an odd prime $p\mid m$, $d=m/p$, $\gcd(a,d)=1$, the
variety $Y$ defined by (2.1) there is a subvariety of $X^{p-1}_m$ of codimension $r=(p-1)/2$
representing $\alpha=\sigma_{p,a}$; more precisely
$\omega_\alpha(Y)\cdot\overline{\omega_\alpha(Y)}=(-1)^r p^{p-2}m^p$.}
\end{itemize}

\emph{From \cite{Aoki00}} (verified verbatim against the open-access copy in the Rikkyo
University repository, doi:10.14992/00009819; \S\ref{sec:closures} records two errata of the
printed text, both likewise verified against the primary):

\begin{itemize}
\item[--] \textbf{\S2 definitions} (verbatim). \emph{If $m$ is divisible by a prime number $p$
and $m>p$, then we have explicitly defined elements of $\mathfrak B_m$:
$\sigma_{p,a}=(a,\,a+\frac mp,\,\dots,\,a+\frac{(p-1)m}p,\,m-pa)$ if $p\ge3$;
$(a,\,a+\frac m2,\,m-2a,\,\frac m2)$ if $p=2$ --- where $a$ is an integer such that
$0<a<\frac mp$. We call $\sigma_{p,a}$ a standard element.} Further (\S2 there): $A_m$ is the
zero-sum sublattice of the free abelian group on $\Z/m\smallsetminus\{0\}$; $u$ (our $\nu$) the
multiplicity-vector map, $\nu(\alpha*\beta)=\nu(\alpha)+\nu(\beta)$; $B_m,S_m,D_m\subset A_m$ the
subgroups generated by $u(\mathfrak B_m),u(\mathfrak S_m),u(\mathfrak D_m)$, with
$\mathfrak D_m\subset\mathfrak S_m\subset\mathfrak B_m$ and $\mathfrak S_m$ the set of $\alpha$
admitting $\alpha*\delta\sim\sigma_1*\cdots*\sigma_k*\delta'$ ($\delta,\delta'\in\mathfrak
D_m$).
\item[--] \textbf{Theorem 0.1.} \emph{Suppose the prime factorization of $m$ is one of the
following forms: (i) $m=2^a3^b5^c7^d$, where $a,b,c,d$ are non-negative integers such that either
$c=0$ or $d=0$; (ii) $m=p^e$ or $2p^e$, $p$ an odd prime. Then the Hodge conjecture is true for
all abelian varieties of Fermat type of degree $m$.} Its proof (\S4 there, p.~185) opens: ``In view of
Corollary~3.2, it suffices to show that the Hodge conjecture for $X^n_m$ is true for all $n$'' ---
the bridge to the Fermat varieties themselves used in Proposition~\ref{prop:w168}; the
verification then reduces, via Lemma~4.1 and Theorem~2.1 there, to the gap generators $\xi_m$ at
$m=12,15,20,21,28$ (the $m=28$ entry being the misprint recorded in \S\ref{sec:closures}).
\item[--] \textbf{Theorem 2.1} (verbatim; proved as Theorem~D of \cite{AokiD}). \emph{For each
divisor $d>1$ of $m$ such that $\mu'(d)=1$, there exists an element $\xi_d$ of $B_d-S_d$ with
the property that the $\Gamma_m$-module $B_m$ is generated by $(m/d)\xi_d$ for all such divisors
$d$ of $m$ and $S_m$. Moreover, we have $2\alpha\in S_m$ for every $\alpha\in B_m$.} ($\mu'$ is
Aoki's modified M\"obius function and $(m/d)$ his level-raising product map; their exact forms
are not needed here.)
\item[--] \textbf{Lemma 4.1} (verbatim). \emph{Let $\alpha,\beta\in B_m$. Then:
(i) if both $\V(\alpha)$ and $\V(\beta)$ are algebraic, then so is $\V(\alpha*\beta)$;
(ii) if $\V(\alpha*\delta)$ is algebraic for some $\delta\in D_m$, then so is $\V(\alpha)$;
(iii) if $\alpha\in S_m$, then $\V(\alpha)$ is algebraic.} No hypothesis on $m$ is imposed;
Proposition~\ref{prop:ident} identifies the implemented lattice with $S_m$
generator-by-generator.
\end{itemize}

\begin{thebibliography}{99}
\bibitem{Aoki87} N.~Aoki, \emph{Some new algebraic cycles on Fermat varieties}, J.~Math.\ Soc.\ Japan
\textbf{39} (1987), 385--396.
\bibitem{AokiD} N.~Aoki, \emph{On some arithmetic problems related to the Hodge cycles on the
Fermat varieties}, Math.\ Ann.\ \textbf{266} (1983), 23--54; Erratum: Math.\ Ann.\ \textbf{267}
(1984), 572.
\bibitem{Aoki00} N.~Aoki, \emph{Some remarks on the Hodge conjecture for abelian varieties of Fermat
type}, Comment.\ Math.\ Univ.\ Sancti Pauli \textbf{49} (2000), 177--194; doi:10.14992/00009819
(open access, Rikkyo University repository).
\bibitem{Aoki02} N.~Aoki, \emph{Hodge cycles on CM abelian varieties of Fermat type}, Comment.\
Math.\ Univ.\ Sancti Pauli \textbf{51} (2002), 99--130; doi:10.14992/00008725.
\bibitem{daSilva} G.~da Silva Jr., \emph{Notes on the Hodge conjecture for Fermat varieties},
Experimental Results \textbf{2} (2021), e22, doi:10.1017/exp.2021.14; arXiv:2101.04739.
\bibitem{BN} R.~Jumagulov, \emph{The Hodge conjecture for Fermat fourfolds of odd degree at
most 199}, companion note, 2026 (arXiv, to appear; the ancillary package is shared).
\bibitem{P1} R.~Jumagulov, \emph{Galois-invariant N\'eron--Severi ranks of Fermat surfaces over
number fields: a Galois module, closed forms, a threshold, and exact tables}, arXiv:2607.17387.
\bibitem{Shioda79} T.~Shioda, \emph{The Hodge conjecture for Fermat varieties}, Math.\ Ann.\ \textbf{245}
(1979), 175--184.
\bibitem{Shioda81} T.~Shioda, \emph{Algebraic cycles on abelian varieties of Fermat type}, Math.\
Ann.\ \textbf{258} (1981/82), 65--80.
\bibitem{SK} T.~Shioda, T.~Katsura, \emph{On Fermat varieties}, T\^ohoku Math.\ J.\ \textbf{31} (1979),
97--115.
\bibitem{Weil49} A.~Weil, \emph{Numbers of solutions of equations in finite fields}, Bull.\ Amer.\
Math.\ Soc.\ \textbf{55} (1949), 497--508.
\bibitem{Yamamoto} K.~Yamamoto, \emph{The gap group of multiplicative relationship of Gauss
sums}, Symp.\ Math.\ \textbf{XV} (1975), 427--440.
\bibitem{Weil52} A.~Weil, \emph{Jacobi sums as ``Gr\"ossencharaktere''}, Trans.\ Amer.\ Math.\ Soc.\
\textbf{73} (1952), 487--495.
\end{thebibliography}
\end{document}
